\documentclass[10pt,a4paper]{amsart}
\usepackage[margin=19mm]{geometry}
\usepackage{amsmath,amssymb,mathtools}
\usepackage[scr=rsfs,cal=euler]{mathalfa}
\usepackage{xfrac}
\usepackage[unicode,hidelinks,bookmarks=false]{hyperref}
\newcommand{\ie}{i.e.\ }
\newcommand{\cf}{cf.\ }
\newcommand{\resp}{respectively\ }
\newcommand{\Subsection}{\subsection}
\providecommand{\nobreakdash}{\nobreak\hbox{-}}
\setlength{\emergencystretch}{2em}
\multlinegap0pt
\usepackage{bm}
\usepackage{bbm}
\usepackage{dsfont}
\usepackage{xspace}
\usepackage{cancel}
\usepackage{mathtools}
\usepackage{amsthm}
\makeatletter
\@ifundefined{theorem}{\newtheorem{theorem}{Theorem}}{}
\makeatother
\title[Learning to Fuse Temporal Proximity Networks: A Case Study i]{Learning to Fuse Temporal Proximity Networks: A Case Study in Chimpanzee Social Interactions}
\author{Yixuan He, Aaron Sandel, David Wipf, Mihai Cucuringu, John Mitani, Gesine Reinert}
\date{Source: arXiv:2502.00302v3 (CC BY 4.0)}
\begin{document}
\maketitle

\noindent\emph{Source and licence.} Learning to Fuse Temporal Proximity Networks: A Case Study in Chimpanzee Social Interactions. Version arXiv:2502.00302v3 (CC BY 4.0), distributed under the Creative Commons Attribution 4.0 International licence, as recorded in the arXiv licence metadata for this exact version and shown on the abstract page https://arxiv.org/abs/2502.00302v3. Full rights notice: licenses/arxiv-2502.00302-CC-BY-4.0.txt.\par\smallskip

\noindent\textit{Editorial scope.} Counted unit: the Theorem whose source label is \texttt{thm:count\_successes\_distribution} in the article (source0.tex lines 369--386, source label \texttt{thm:count\_successes\_distribution}), delivered together with the article's own complete original proof of it (source0.tex lines 388--409, one \texttt{proof} environment of 22 lines). No mathematics is written, repaired or re-derived: statement and proof are the article's own text line for line. Required context retained uncounted: none. Every definition, display and label of those lines is retained unchanged. Omitted from this package and not redistributed: 1--368, 387--387, 410--795. Nothing inside a retained excerpt is omitted or altered: every retained body is delivered exactly as the article's lines are written, so no omission receipt of any kind accompanies this package. Citations: the retained excerpts carry no citation command at all, so no bibliography entry of the article is transcribed and no reference is redistributed. Reference adaptation: none was needed and none was made; every reference inside the delivered unit resolves inside it, so no unresolved reference or citation is produced. Printed slips kept literal: no printed word, symbol, hypothesis or number of the counted statement or of its proof was altered. Layout and class: the article's own class is \texttt{article} with packages including microtype, graphicx, caption, subcaption, bbm, booktabs, multirow, enumitem, hyperref, arxiv\_style, amsmath, amssymb, mathtools, amsthm; the wrapper is the lane's pinned \texttt{amsart} editorial wrapper at 10pt on A4, which restates the theorem-like environments the retained text needs and adds the article's own macro definitions that the retained text needs (none), with extra wrapper packages (bm, bbm, dsfont, xspace, cancel, mathtools). The delivered numbering is the unit's own. Assets: no retained span contains a figure environment, a graphics-inclusion command, a PDF-page inclusion, a table or a TikZ picture; no image file is copied into this package, the package declares no asset dependency and no image file is redistributed with it. Licence: version arXiv:2502.00302v3 of this article is distributed under the Creative Commons Attribution 4.0 International licence, as recorded in the arXiv licence metadata for this exact version and shown on the abstract page https://arxiv.org/abs/2502.00302v3. A full rights notice is delivered at licenses/arxiv-2502.00302-CC-BY-4.0.txt. Characters: every retained excerpt is pure ASCII, so the delivered engine, XeLaTeX with its default Latin Modern text font, reports no missing character.
\begin{theorem}
\label{thm:count_successes_distribution}
[Useful for Count Similarity] For a 
series of independent Bernoulli trials $\{B_t\}$ of length $T\;(T\geq 2)$, $t\in \{1, \dots, T\},$ with success rate $p_t$ for $B_t,$ denote $C_t \in \{0, \dots, t\}$ as the total number of successes in time steps $\{ 1, \dots, t\}$. 
For $t\in\{1, \dots, T\}$, the probability distribution of $C_t$ satisfies
\begin{align}
\mathbb{P}(C_t=0) &=\prod_{s=1}^t(1-p_s), \;\;
 \mathbb{P}(C_t=t) =\prod_{s=1}^t p_s,
 \end{align}
 and for $L\in\{1, \dots, t-1\}$ if $t\geq 2,$
 \begin{align}
 \mathbb{P}(C_t=L) &=p_t\cdot\mathbb{P}(C_{t-1}=L-1) + (1-p_t)\cdot \mathbb{P}(C_{t-1}=L);
    \end{align}
    for $ L\geq t+1,$
    \begin{align}
    \mathbb{P}(C_t=L)&=0.
\end{align}
\end{theorem}

\begin{proof}

We first prove the behavior at the boundaries.

Since the total number of successes is bounded by the number of trials, we always have $C_t\leq t.$ Therefore, $\mathbb{P}(C_t=L)=0$ whenever $ L\geq t+1.$ The only chance of having pure successes is to never fail the trial, and hence $\mathbb{P}(C_t=t) =\prod_{s=1}^t p_s, \;\; t\in\{1, \dots, T\},$ due to the independence between the Bernoulli random variables. Similarly, the only chance of having zero successes is to fail the trial every time, and hence $\mathbb{P}(C_t=0) =\prod_{s=1}^t(1-p_s), \;\; t\in\{1, \dots, T\}$.

In order to compute $\mathbb{P}(C_t=L)$ for $t\in\{2, \dots, T\},\;L\in\{1, \dots, t-1\},$ we first analyze the role of this time step $t$. There are two possibilities for time step $t$: either it is a success, or it is a ``failure" point. Note that we concentrate on the time series from the start (time step $1$) till time step $t,$ and we impose no constraints on future time steps. For the first situation, we require that we have $L-1$ successes before $t;$ while for the second situation, we require $L$ successes at time $t-1.$ Therefore, 
\begin{align*}
    \mathbb{P}(C_t=L) &=p_t\cdot\mathbb{P}(C_{t-1}=L-1)+ (1-p_t)\cdot \mathbb{P}(C_{t-1}=L), \\
    & \qquad \qquad \;\;t\in\{2, \dots, T\},\;L\in\{1, \dots, t-1\}.
\end{align*}
Combing the above, we have that 
\begin{align*}
\begin{split}
    \mathbb{P}(C_t=t) &=\prod_{s=1}^t p_s, \;\; t\in\{1, \dots, T\},\\
    \mathbb{P}(C_t=0) &=\prod_{s=1}^t(1-p_s), \;\; t\in\{1, \dots, T\},\\
    \mathbb{P}(C_t=L) &=p_t\cdot\mathbb{P}(C_{t-1}=L-1)+ (1-p_t)\cdot \mathbb{P}(C_{t-1}=L),\\
    & \qquad \qquad  \;\;t\in\{2, \dots, T\},\;L\in\{1, \dots, t-1\},\\
    \mathbb{P}(C_t=L)&=0,\;\;t\in\{1, \dots, T-1\}, \; L\in\{t+1, \dots, T\}.
\end{split}
\end{align*}
\end{proof}

\end{document}
